\begin{textsample}{NEG Dim 6 – human – Score -2.00 – es/es\_ef\_anos\_finais\_linguagens\_ingles\_1\_p8.txt}  \label{ex:f6_neg_016}
A partir do entendimento dos Multiletramentos como usos e práticas sociais da linguagem relacionados aos seus contextos socio‐históricos de produção, as práticas de \textbf{letramento} devem ser concebidas como práticas de \textbf{leitura}, de escrita e de oralidade em diferentes semioses dentro de uma perspectiva social e cultural.
As diferentes semioses consideram o texto para além do escrito, incluindo elementos de representação espacial, corporal, sonora e visual, como os hipertextos, muitas vezes multimidiáticos, dados os meios de circulação que hoje fazem parte do universo “ conectado ” dos estudantes.
Em sua essência, os Multiletramentos podem ser compreendidos como práticas sociais de linguagem situadas em diversos contextos, como pode ser \textbf{observado} nos Eixos de Práticas de Linguagem aqui sugeridos.

% matched lemmas: leitura, letramento, observar
\end{textsample}
